\section{Conclusión}
A partir de la metodología cualitativa aplicada en este proyecto, se comprobó que la integración de herramientas de control de versiones, metodologías ágiles, arquitecturas de software y lenguajes robustos es esencial para el éxito en el desarrollo de sistemas. El análisis documental y las evidencias prácticas permitieron identificar como pilares fundamentales a Git y GitHub, la Guía de Scrum 2020, los patrones de diseño, el patrón Modelo–Vista–Controlador (MVC) y el lenguaje Java, cada uno aportando ventajas específicas para la organización, la colaboración y la escalabilidad de los proyectos. Los resultados obtenidos demuestran que el uso de Git y GitHub garantizó la trazabilidad de cambios y el trabajo colaborativo; Scrum 2020 fortaleció la planificación, la comunicación y la entrega incremental de funcionalidades; los patrones de diseño y la arquitectura MVC facilitaron la estructura clara del sistema; y Java proporcionó robustez, portabilidad y un ecosistema amplio para el desarrollo. En conjunto, estos hallazgos evidencian que la combinación de una base teórica sólida con la práctica constante en los tres proyectos realizados durante el año permitió obtener productos funcionales y alineados con las necesidades planteadas, confirmando la eficacia de esta integración de tecnologías y metodologías en entornos académicos y profesionales.